\section{La paradoja de las tareas largas: entre más pueda correr el Agent, más se cansa el usuario}

La observación más cercana a la experiencia cotidiana en este episodio es la del context switching. Shang Yanyi (尚晏仪) dice que, después de usar mucho Deep Research y coding agents, siente una carga muy fuerte: se lanza un prompt completo y el Agent corre durante mucho tiempo; si la capacidad es mala, ella no lo usa; si la capacidad es buena, tampoco necesita ver la salida intermedia. Entonces no sabe qué hacer mientras espera, así que cambia a otra tarea y regresa unos minutos después a revisar, lo cual la cansa aún más.

Esta observación es muy importante porque refuta buena parte de la narrativa por defecto de muchos productos de Agent. Hoy varios productos usan «poder hacer una tarea de una hora» como argumento de venta, pero Gu Yu (谷雨) señala que si se pudiera lograr la misma calidad en 30 segundos, nadie querría esperar una hora. La capacidad de sostener tareas largas sin duda tiene valor, pero es una capacidad del sistema, no la experiencia del usuario en sí misma. Presentar la lentitud como autonomía es una narrativa de producto muy peligrosa.

\begin{importantbox}{La capacidad de tareas largas no es una experiencia de espera larga}
Lo que el usuario quiere es que la tarea se complete de forma confiable, no verse obligado a administrar una caja negra que corre durante mucho tiempo. Que el Agent pueda correr mucho tiempo indica que quizás posee capacidad de ejecución autónoma; pero el valor del producto viene de menos espera, menos revisión, menos cambio de contexto cognitivo y una mayor calidad de resultado.
\end{importantbox}

La conversación también divide este problema en capa de producto y capa de capacidad. La capa de producto puede mejorar la sensación mediante estados intermedios, retroalimentación de progreso, mecanismos de interrupción y puntos de confirmación clave. Por ejemplo, algunos coding agents le informan continuamente al usuario qué están haciendo, mientras que otros sistemas pasan varios minutos sin dar salida alguna y eso genera ansiedad en el usuario. Pero la capa de capacidad es más fundamental: que el Agent busque información con lentitud, use herramientas sin soltura, siga rutas de decisión enredadas y abra repetidamente páginas web inútiles indica que todavía no posee esa capacidad de juzgar rápidamente el valor de la información que tiene un humano experimentado.

Esto también explica por qué «escribir un prompt para que el Agent corra seis horas» solo le sirve a una minoría de usuarios. Las tareas reales suelen necesitar aclaraciones e iteraciones de ida y vuelta. Ni siquiera el propio usuario puede necesariamente escribir de una vez un requerimiento completo. Si un Agent toma un requerimiento ambiguo y corre durante mucho tiempo, es muy probable que al final produzca una gran cantidad de trabajo que parece esforzado pero que en realidad se desvió. La autonomía real no debería ser «correr solo durante mucho tiempo», sino «saber cuándo avanzar, cuándo detenerse y cuándo pedirle un juicio al usuario».

\begin{knowledgebox}{Cuatro puntos de diseño para la experiencia de tareas largas}
\begin{itemize}
  \item Compresión de estado: decirle al usuario qué problema se está resolviendo en este momento, en lugar de volcarle los registros de bajo nivel.
  \item Puntos de verificación: pedir confirmación en las bifurcaciones de alto riesgo, en lugar de entregar una sorpresa de una sola vez al final.
  \item Interrumpibilidad: permitir que el usuario cambie de rumbo, pause o agregue información.
  \item Resumen de diferencias: al terminar, explicar qué se hizo, qué se cambió y qué sigue siendo incierto.
\end{itemize}
\end{knowledgebox}

\subsection{Resumen de la sección}

Una nueva contradicción de los productos de Agent es: entre más se parezca la capacidad a «ejecución autónoma de larga duración», más puede empujar al usuario hacia revisiones frecuentes y cambios de contexto. La capa de producto puede aliviar esto mediante el diseño de retroalimentación, pero lo fundamental sigue siendo la capacidad del modelo en tool use, decision making y juicio de tareas. A largo plazo, el usuario no debería entrenarse para adaptarse a un Agent ineficiente; el Agent debería aprender a reducir la carga cognitiva del usuario.

